% !TeX root = ../colloquium-2.tex

% 1-я теорема Вейерштрасса.
% 2-я теорема Вейерштрасса.
% Теорема Коши о непрерывной функции.

\phantomsection\label{sec_Definition02}
\begin{definition}[
    1-я теорема Вейерштрасса.
    2-я теорема Вейерштрасса.
    Теорема Коши о непрерывной функции.
]
    \subsection*{1-я теорема Вейерштрасса.}
    \begin{note}[purple]
        Пусть функция $f$ определена и непрерывна на отрезке $[a, b]$. 
        Тогда она ограничена на этом отрезке.
    \end{note}

    \subsection*{2-я теорема Вейерштрасса.}
    \begin{note}[purple]
        Пусть функция $f$ определена и непрерывна на отрезке $[a, b]$. 
        Тогда существуют такие точки $x_1, \; x_2 \in [a, b]$, что 
        $$
            f(x_1) = m = \inf_{a \leq x \leq b} f(x), \ f(x_2) = M = \sup_{a \leq x \leq b} f(x)
        $$
    \end{note}

    \subsection*{Теорема Больцано-Коши о непрерывной функции.}
    \begin{note}[purple]
        Пусть функция $f$ определена и непрерывна на отрезке 
        $[a, b]$, $m = \displaystyle{\inf_{a \leq x \leq b} f(x)}$,
        $M = \displaystyle{\sup_{a \leq x \leq b} f(x)}$. \\

        Тогда для любого числа $C \in [m, \; M]$ существует такая точка 
        $c \in [a, \; b]$, что $f(c) = C$. 
    \end{note}
\end{definition}